\title{Spectral Radius Minimization for Signed Squares of Cycles}
\ifdefined\TargetAAnonymous
  \author{Anonymous}
\else
  \author{[AUTHOR NAME]}
\fi
\date{}
\maketitle

\ifdefined\TargetAAnonymous\else
\begin{center}
\small
\mbox{[AFFILIATION]}\\
\mbox{[DEPARTMENT]}, \mbox{[INSTITUTION]}\\
\mbox{[CITY]}, \mbox{[COUNTRY]}\\[3pt]
Corresponding author: [CORRESPONDING AUTHOR]\\
Email: [EMAIL]\quad ORCID: [ORCID]
\end{center}
\fi

\begin{abstract}
For even \(n\geq 8\), let \(C_n^2=C_n(1,2)\) be the square of the cycle,
with its distance-one and distance-two edges independently assigned signs,
with no periodicity assumption on the signing, and let \(m_n\) be the minimum
signed adjacency spectral radius.  A
distinguished twisted signing attains \(\rho_-(n)\), where
\(\rho_-(n)^2=4+2\cos(2\pi/n)+2\cos(4\pi/n)\), at every even order.  We prove
that equality holds exactly for
\(n\in\{8,10,12,14,16,18,20,22,24,26,28,30,34,36,38,42,44,46\}\); equivalently,
strict failure occurs precisely at \(32\), at \(40\), and at every even
\(n\geq48\).  Structurally, a period-eight reference phase has squared
spectral edge \(\eta<8\), while its six-gap phase slip is the unique
least-edge abnormal positive single gap and has an isolated rank-two squared
edge \(c_6\in(\eta,8)\).  Charge-compatible separated copies of this slip,
together with a discrete IMS estimate, prove eventual failure, while exact
finite verification shows that continuous failure begins precisely at order
\(48\).  The finite portions are reduced to explicitly finite objects and
independently verified using exact arithmetic, so floating-point evidence
makes no spectral decision.
\end{abstract}

\noindent\textbf{Keywords:} signed graph; spectral radius; cycle square;
switching equivalence; phase slip; computer-assisted proof

\medskip
\noindent\textbf{MSC 2020:} 05C50 (primary); 05C22, 68V05 (secondary)
